% Hafta 5 — SEI CERT Oracle Java: aynı yapı, farklı dil
% Derleme: py -3.12 tools/latex/derle.py docs/week-5/assets/h05-11-cert-java.tex
\documentclass[tikz,border=8pt]{standalone}
\usepackage{cen429-cizim}
\begin{document}
\begin{tikzpicture}
  \node[baslik] (b) at (0,0) {SEI CERT Oracle Java: aynı yapı, farklı dil};
  \node[altbaslik] at (b.south west) {yönetilen dil bellek hatasını çözer, mantık hatasını çözmez};
  \node[morkutu=34mm, anchor=north west] (n1) at (0,-2.4)
    {\kb{Kural kimliği}\\IDS00-J, SER01-J\\alan + sıra};
  \node[kotukutu=34mm, right=9mm of n1] (n2)
    {\kb{Hatalı kod}\\yönetilen dilde de\\aynı hata};
  \node[iyikutu=34mm, right=9mm of n2] (n3)
    {\kb{Uyumlu kod}\\API'nin güvenli\\kullanımı};
  \node[uyarikutu=32mm, right=9mm of n3] (n4)
    {\kb{Risk}\\ciddiyet · olasılık\\· düzeltme};
  \draw[ok] (n1.east) -- (n2.west);
  \draw[ok] (n2.east) -- (n3.west);
  \draw[ok] (n3.east) -- (n4.west);
  \node[fit=(n1)(n2)(n3)(n4), inner sep=0pt] (grp) {};
  \node[dip, text width=155mm, anchor=north west] at ($(grp.south west)+(0,-8mm)$)
    {IDS (girdi doğrulama), SER (serileştirme), FIO (dosya), MSC (çeşitli) --- enjeksiyon konuları IDS'tedir.};
\end{tikzpicture}
\end{document}
